% Part: applied-modal-logics
% Chapter: temporal-logic
% Section: possible-histories

\documentclass[../../../include/open-logic-section]{subfiles}

\begin{document}

\olfileid[id]{aml}{tl}{poss}

\olsection{Riwayat yang Mungkin}

Model relasional logika temporal yang telah kita gunakan sangat luwes karena
kita tidak perlu mengenakan pembatasan apa pun pada relasi aksesibilitas. Ini
berarti model temporal dapat bercabang pada masa lampau maupun masa depan,
tetapi kita mungkin ingin mempertimbangkan konsepsi percabangan yang lebih
``modal,'' tempat kita memandang barisan peristiwa sebagai riwayat yang
mungkin. Hal ini tidak mesti mengharuskan kita mengubah bahasa, meskipun kita
juga dapat menambahkan operator modal ``biasa'' kita, $\Box$ dan~$\Diamond$,
serta mempertimbangkan penambahan relasi aksesibilitas epistemik untuk
merepresentasikan perubahan pengetahuan agen dari waktu ke waktu.

\begin{defn}
  Suatu \emph{model riwayat mungkin} bagi bahasa temporal ialah tripel
  $\mModel{M}=\tuple{T,C,V}$, dengan
  \begin{enumerate}
    \item $T$ merupakan himpunan tak kosong yang ditafsirkan sebagai
      keadaan-keadaan dalam waktu.
    \item $C$ merupakan himpunan lintasan komputasi, atau \emph{riwayat yang
      mungkin}, dari suatu sistem. Dengan kata lain, $C$ merupakan himpunan
      barisan~$\sigma$ dari keadaan $s_1$, $s_2$,~$s_3$, \dots, dengan setiap
      $s_i\in T$.
    \item $V$ merupakan fungsi yang menetapkan kepada setiap !!{propositional
      variable}~$p$ suatu himpunan~$V(p)$ yang terdiri atas titik-titik waktu.
  \end{enumerate}
  Untuk menyederhanakan pembahasan, pada umumnya kita juga mengandaikan bahwa
  jika suatu riwayat termasuk dalam~$C$, demikian pula semua sufiksnya.
  Sebagai contoh, jika $s_1$, $s_2$,~$s_3$ merupakan barisan dalam~$C$,
  demikian pula $s_2$, $s_3$ dan~$s_3$.
  Selain itu, jika dua keadaan $s_i$ dan~$s_j$ muncul dalam suatu
  barisan~$\sigma$, kita menulis $s_i\prec_\sigma s_j$ jika $i<j$. Jika
  $t\in V(p)$, kita mengatakan bahwa $p$ \emph{benar pada}~$t$.
\end{defn}

Satu perubahan yang relevan ialah bahwa ketika kita mengevaluasi kebenaran
!!a{formula} pada titik waktu~$t$ dalam model~$\mModel M$, kita melakukannya
relatif terhadap suatu riwayat~$\sigma$ yang memuat $t$ sebagai keadaan.
Namun, kita tidak perlu mengubah semantik !!{propositional variable}s atau
konektif fungsional-kebenaran. Semuanya tetap persis seperti dalam
\olref[sem]{defn:tmodels}, sebab tidak satu pun di antaranya mengacu
pada~$\sigma$. Akan tetapi, sekarang kita mendefinisikan ulang operator masa
depan~$\Ftemp$ dan menambahkan operator~$\Diamond$ berkenaan dengan
riwayat-riwayat tersebut.

\begin{defn}\ollabel{defn:phmodels}
  \emph{Kebenaran !!a{formula}~$!A$ pada~$t,\sigma$} dalam~$\mModel M$,
  dengan simbol $\mSat{M}{!A}[t,\sigma]$:
  \begin{enumerate}
  \item\ollabel{defn:sub:phmodels-f} \indcase{!A}{\Ftemp
    !B}{$\mSat{M}{\indfrm}[t,\sigma]$ jika dan hanya jika
    $\mSat{M}{!B}[t',\sigma]$ untuk suatu $t'\in T$ dengan
    $t\prec_\sigma t'$.}
  \item\ollabel{defn:sub:phmodels-diamond} \indcase{!A}{\Diamond
    !B}{$\mSat{M}{\indfrm}[t,\sigma]$ jika dan hanya jika
    $\mSat{M}{!B}[t,\sigma']$ untuk suatu $\sigma'\in C$ yang memuat~$t$.}
  \end{enumerate}
\end{defn}

Operator temporal dan modal lainnya dapat didefinisikan dengan cara serupa.
Namun, sekarang kita dapat merepresentasikan pernyataan yang menggabungkan
kala dan modalitas. Sebagai contoh, kita dapat menyimbolkan ``$p$ tidak akan
terjadi, tetapi mungkin saja terjadi'' dengan formula
$\lnot\Ftemp p\land\Diamond\Ftemp p$. Formula ini berlaku pada suatu titik
dan riwayat tempat $p$ tidak menjadi benar pada keadaan penerus, tetapi
terdapat riwayat alternatif tempat $p$ akan menjadi benar.

\end{document}
